% ---------------------------------------------------------------- discussion
\subsection{Discussion}
\label{sec:os-discussion}

Adding a label-space shift to a geographic one does not give adaptation more to work with.
The cross-building penalty is $1.6$ points entering $\text{S}_2$ and $36.0$ entering
$\text{S}_1$ once classes are withheld, a wider asymmetry than the closed-set
$0.6$ and $21.3$, and in neither direction does any of the four methods improve on a frozen
source model once H is read beside its components. The combined shift is harder, not
richer.

What does help is replacing the rejection score, and only modestly: between $1.17$ and
$1.58$ times chance at a matched rejection rate, against a maximum softmax baseline that
sits exactly at chance. This is a consistent reading alongside
Section~\ref{sec:res-mech}, where DIGA's prototype classifier was the one component that
helped without touching normalisation statistics, but we did not isolate that component in
the open-set setting and the two results are not a replication of one another. What they
have in common is the direction of the leverage: in both, the useful part was a distance
in representation space, and the parts that manipulate source statistics were harmful or
inert.

Three things we would want a reader to take from this section are not about our system.

Metrics first. UCDR has a trivial maximiser and is conventionally paired with OS$^{*}$ to
remove it. That pairing is insufficient, because H is itself partly maximised by uniform
degradation whenever unknown detection is the weaker component, which it usually is.
PETAL records the best H in one direction while classifying the known activities half as
well as doing nothing. Open-set results should be reported with both components visible.

Thresholds second. A threshold set to reject $1\%$ of source validation rejects between
$7\%$ and $63\%$ of a target stream, depending only on which score carries it. Any
open-set method that calibrates on source inherits this, and any that calibrates on target
has left the test-time setting. We think the rejection rate, which is observable without
labels, belongs in open-set reporting next to H. How to calibrate a threshold that holds
across domains rather than across resamples of one is open, and is the direction we would
pursue next.

Chance baselines third. The two results in this section that are most likely to generalise
are both negative, and neither is visible without a chance baseline: four standard
confidence scores are at or below chance as novelty detectors, and a learned representation
carries no more novelty signal than the untrained input. Neither would appear as a problem
in a table of H-scores alone.

The limits are the same as the closed-set study's. Two buildings, one dataset, one choice
of withheld activities. The asymmetry we report is a property of this pair, and the
specific values of $1.6$ and $36.0$ belong to it. What we would expect to transfer is the
procedure: measure the penalty per cell against a target-trained model at matched sample
size before running any method, report rejection rates beside H, and compute a chance
baseline for any novelty score before believing it.
